\documentclass[10pt,a4paper]{amsart}
\usepackage[margin=19mm]{geometry}
\usepackage{amsmath,amssymb,mathtools}
\usepackage[scr=rsfs,cal=euler]{mathalfa}
\usepackage{xfrac}
\usepackage[unicode,hidelinks,bookmarks=false]{hyperref}
\newcommand{\ie}{i.e.\ }
\newcommand{\cf}{cf.\ }
\newcommand{\resp}{respectively\ }
\newcommand{\Subsection}{\subsection}
\providecommand{\nobreakdash}{\nobreak\hbox{-}}
\setlength{\emergencystretch}{2em}
\multlinegap0pt
\usepackage{amssymb}
\usepackage{mathtools}
\usepackage[shortlabels]{enumitem}
\usepackage{xcolor}
\newtheorem{lemma}{Lemma}
\newtheorem{theorem}{Theorem}
\usepackage{cleveref}
\crefname{lemma}{Lemma}{Lemmas}
\crefname{theorem}{Theorem}{Theorems}
\newcommand{\Frac}{\operatorname{Frac}}
\newcommand{\Span}{\operatorname{Span}}
\newcommand{\Spec}{\operatorname{Spec}}
\newcommand{\fm}{\mathfrak m}
\title{Polynomial extensions do not preserve the strong finite type property}
\author{Viet-Hoang Tran, Phan Thanh Toan, Thieu N. Vo, Tan M. Nguyen}
\date{Source: arXiv:2608.26044v1 (CC BY 4.0)}
\begin{document}
\maketitle

\noindent\emph{Source and licence.} Polynomial extensions do not preserve the strong finite type property. Version arXiv:2608.26044v1 (CC BY 4.0), distributed by arXiv under the Creative Commons Attribution 4.0 International licence, http://creativecommons.org/licenses/by/4.0/, as recorded in the arXiv licence metadata for this exact version and shown on the abstract page https://arxiv.org/abs/2608.26044v1. Full licence notice: \texttt{licenses/arxiv-2608.26044-CC-BY-4.0.txt}.\par\smallskip

\noindent\textit{Editorial scope.} Counted unit: the article's theorem thm:equicharacteristic-zero (source0.tex lines 728--732). The article's complete original proof of it is delivered with it (source0.tex lines 734--907, one \texttt{proof} environment of 174 lines). No mathematics is written, repaired or re-derived: the statement and the proof are the article's own lines, in the article's own notation, and no retained line is substituted or re-flowed.  Retained uncounted as the source-local context the proof needs: lines 335--340; lines 473--479.  Nothing was omitted from any retained span, so the package carries no omission record.  No retained line contains a citation, so the wrapper carries no bibliography and no bibliography entry.  No retained line contains a figure environment, an image-inclusion command or drawing code, so no image is delivered and the package declares no asset dependency.  Editorial formatting only, in addition: an AMS \texttt{amsart} preamble that restates the article's own declarations and macros used by the retained lines, namely the environments \texttt{lemma}, \texttt{theorem} on separate counters, together with the standard packages those lines need.  The retained environments print in this document as Lemma 1, Theorem 1 in order of appearance, whereas the article prints the same items with the numbering of its own full document.  The complete native source of the article as distributed in the arXiv e-print for this exact version is bundled unchanged as \texttt{sources/208617/source0.tex}.
\begin{lemma}\label{lem:formal-expansion}
Every $r\in R$ has a unique formal $c$-adic expansion
\begin{equation}\label{eq:formal-expansion}
  r=\sum_{n\geq0}c^nr_n\in K[[c]],\qquad r_n\in V_n.
\end{equation}
\end{lemma}

\begin{lemma}\label{lem:division-char2}
Let $d\geq1$, and let $h\in R[X]$ have degree at most $d$ and satisfy
$h(u)=0$ in $K(c)$.  Write $h=(X-u)\widetilde h$ in $K(c)[X]$.  Then
$\widetilde h\in K[c]_{(c)}[X]$, and, for every $n\geq0$, the coefficient
of $c^n$ in each $X$-coefficient of $\widetilde h$ belongs to
$V_{n+d-1}$.
\end{lemma}

\begin{theorem}\label{thm:equicharacteristic-zero}
There is a one-dimensional local SFT domain $R_0$ containing $\mathbb Q$
such that $R_0[X]$ is not SFT.  More precisely, $R_0[X]$ has a non-SFT
upper to zero.
\end{theorem}

\begin{proof}
Let $k=\overline{\mathbb Q}$, a countable field containing every root of
unity.  Let $u$ be transcendental over $k$ and put $K=k(u)$.  The field $K$
is countable, so enumerate
$K^\times=\{\xi_1,\xi_2,\ldots\}$ with $\xi_1=1$.
Use the labelled alphabet
\[
  \{a_j:j\geq0\}\amalg\{x_i:i\geq1\},
  \qquad a_j\mapsto u^{2^j},\quad x_i\mapsto\xi_i,
\]
with weights $w(a_j)=1$ and $w(x_i)=i$.  Labels are retained even when
their images in $K$ agree.  For $n\geq0$, define
\begin{equation}\label{eq:V-char0}
 V_n:=\Span_k\left\{
  \text{images of words of total weight at most }2^n-1
 \right\}\subseteq K,
\end{equation}
where the empty word is allowed.

These spaces are increasing and exhaustive, with $V_0=k$.  Concatenation
of words and the elementary budget inequalities give
\begin{equation}\label{eq:filtration-char0}
 V_nV_m\subseteq V_{n+m},\qquad
 V_iV_jV_\ell\subseteq V_{i+j+\ell-1}\quad(i,j,\ell\geq1).
\end{equation}
Indeed, if $B_r=2^r-1$, then
$B_n+B_m\leq B_{n+m}$ and
$B_i+B_j+B_\ell\leq B_{i+j+\ell-1}$ for positive indices.  We also have
\begin{equation}\label{eq:powers-char0}
 u^{2^j}\in V_1\quad(j\geq0),\qquad u^t\in V_t\quad(t\geq0).
\end{equation}

We need the escape property
\begin{equation}\label{eq:escape-char0}
  k[u]\nsubseteq V_E\qquad(E\geq0).
\end{equation}
Fix $E$ and set $B=2^E-1$.  A word of weight at most $B$ has image
$h(u)u^s$, where $h$ belongs to a finite set of products of
$\xi_1,\ldots,\xi_B$, and $s$ is a sum of at most $B$ powers of $2$.
Let $S_B$ be the set of all such $s$.  Then
\[
  \#(S_B\cap[0,L])=O_B((\log L)^B).
\]
After clearing one common denominator for the finite set of possible $h$,
there are a nonzero $\Delta\in k[u]$ and a finite set $F\subseteq\mathbb N$
such that
\[
  \operatorname{supp}(\Delta f)\subseteq F+S_B\qquad(f\in V_E).
\]
If $k[u]\subseteq V_E$, then, for any
$r\in\operatorname{supp}(\Delta)$, this would force
$r+\mathbb N\subseteq F+S_B$.  The left side has linear growth, whereas
the right side has $O_B((\log L)^B)$ elements up to $L$.  This proves
\eqref{eq:escape-char0}.

Let $c$ be transcendental over $K$ and set
\[
 A_0=\bigoplus_{n\geq0}c^nV_n\subseteq K[c],\qquad
 \fm_0=(A_0)_+,
 \qquad R_0=(A_0)_{\fm_0},\qquad M_0=\fm_0R_0.
\]
By \eqref{eq:filtration-char0}, $A_0$ is a graded domain and
$A_0/\fm_0\cong k$, so $R_0$ is local with maximal ideal $M_0$.  The proof
of \cref{lem:formal-expansion} uses only multiplicativity and therefore
shows that every $r\in R_0$ has a unique expansion
\begin{equation}\label{eq:formal-char0}
  r=\sum_{n\geq0}c^nr_n\in K[[c]],\qquad r_n\in V_n.
\end{equation}
The second inclusion in \eqref{eq:filtration-char0} gives
\begin{equation}\label{eq:cube-char0}
  z^3\in cR_0\qquad(z\in M_0),
\end{equation}
because every coefficient of $b^3/c$, for $b\in\fm_0$, is a sum of
triple products whose filtration indices drop by one; localization gives
the assertion for $M_0$.

We next show that every $0\neq a\in M_0$ divides a power of $c$.  Write
$a=b/t$ with $b\in\fm_0$ and $t\in A_0\setminus\fm_0$.  Since $t$ is a
unit, $aR_0=bR_0$, so it suffices to work with
\[
  b=\sum_{j=\nu}^{L}c^jb_j\in\fm_0,
  \qquad b_j\in V_j,\quad b_\nu\neq0.
\]
Choose $i_0$ with $\xi_{i_0}=b_\nu^{-1}$ and put $w_0=i_0$.
Choose $q$ so large that
\[
  q(2^L-1+w_0)\leq2^q-1.
\]
Since $T^q-1$ has exactly $q$ distinct roots in $k$, put
$\mu_q=\{\zeta\in k:\zeta^q=1\}$ and form
\[
  \mathcal N_q(b):=\prod_{\zeta\in\mu_q}b(\zeta c).
\]
After factoring the lowest term, invariance under $c\mapsto\eta c$ for
$\eta\in\mu_q$ gives
\begin{equation}\label{eq:norm-char0}
 \mathcal N_q(b)=\varepsilon c^{\nu q}b_\nu^q
  \left(1+\sum_{r\geq1}\theta_r c^{rq}\right),
  \qquad \varepsilon\in k^\times.
\end{equation}
Only finitely many $\theta_r$ are nonzero.
Each ratio $b_j/b_\nu$ is a $k$-linear combination of words of weight at
most $2^L-1+w_0$.  Hence every product of at most $q$ such ratios has weight
at most $q(2^L-1+w_0)\leq2^q-1$.  Thus
$\theta_r\in V_q\subseteq V_{rq}$.  The same bound applied to the $q$-fold
product of the label for $b_\nu^{-1}$ gives $b_\nu^{-q}\in V_q$.
Therefore
\[
 \lambda:=c^qb_\nu^{-q}\in A_0,
 \qquad \omega:=\sum_{r\geq1}\theta_rc^{rq}\in\fm_0.
\]
Every $b(\zeta c)$ lies in $A_0$, since $\zeta\in k$, and the factor for
$\zeta=1$ is $b$.  Thus $\mathcal N_q(b)\in bA_0$, and
\eqref{eq:norm-char0} yields
\[
 c^{\nu q+q}=\varepsilon^{-1}(1+\omega)^{-1}
   \lambda\mathcal N_q(b)\in bR_0=aR_0.
\]

Let $I$ be a nonzero proper ideal of $R_0$ and choose $0\neq a\in I$.
Since $I\subseteq M_0$, for some $D$ the preceding paragraph gives
$c^D\in aR_0$.  By
\eqref{eq:cube-char0}, every $z\in I$ satisfies
$z^{3D}\in c^DR_0\subseteq aR_0$.  Thus $I$ is SFT with data
$(I,(a),3D)$, and $R_0$ is SFT.  The same two facts show that every nonzero
prime contains $c$ and then $M_0$; hence
\[
  \Spec R_0=\{(0),M_0\},\qquad \dim R_0=1.
\]
In particular, $R_0$ is equicharacteristic zero, with residue field $k$.

Exhaustivity gives $\Frac(R_0)=K(c)$.  Define
\[
 Q_0:=\ker\bigl(R_0[X]\longrightarrow K(c),\ X\longmapsto u\bigr).
\]
This is a prime with $Q_0\cap R_0=(0)$, and it is nonzero because, for
$m=2^n$,
\begin{equation}\label{eq:gn-char0}
  g_n:=c(X^m-u^m)\in Q_0;
\end{equation}
indeed, $u^m\in V_1$.  Thus $Q_0$ is an upper to zero.

Suppose that $Q_0$ had SFT data $(Q_0,J,e)$, with
$J=(h_1,\ldots,h_t)$.  The ideal $J$ is nonzero because $R_0[X]$ is a
domain and $g_0\neq0$.  Discard zero generators, put
$d=\max_i\deg h_i\geq1$, and fix an integer $N\geq e$.  The proof of
\cref{lem:division-char2}, using \eqref{eq:powers-char0}, writes
$h_i=(X-u)\widetilde h_i$ so that every $c^\ell$-coefficient of every
$X$-coefficient of $\widetilde h_i$ lies in $V_{\ell+d-1}$.
Since $g_n^N\in J$, cancellation of $X-u$ gives
\begin{equation}\label{eq:cancelled-char0}
 c^N\frac{(X^m-u^m)^N}{X-u}
   =\sum_{i=1}^ta_{i,n}\widetilde h_i,
  \qquad a_{i,n}\in R_0[X].
\end{equation}
Use \eqref{eq:formal-char0} and compare the coefficient of $c^N$.
The convolution is finite, and a term with $c$-degrees $s$ and $\ell$ on
the right belongs to $V_sV_{\ell+d-1}\subseteq V_{N+d}$.  Consequently
every coefficient of
\[
  G_{m,N}:=\frac{(X^m-u^m)^N}{X-u}
\]
lies in the fixed space $V_{N+d}$.

For $0\leq t<mN$, set $s=\lfloor t/m\rfloor+1$.  Direct division gives
\begin{equation}\label{eq:coefficients-char0}
 [X^t]G_{m,N}
  =(-1)^{N-s}\binom{N-1}{s-1}u^{mN-1-t}.
\end{equation}
The scalar is nonzero because $k$ has characteristic zero.  Hence
$1,u,\ldots,u^{N2^n-1}\in V_{N+d}$ for every $n$.  Letting $n$ grow
contradicts \eqref{eq:escape-char0}.  Thus $Q_0$ is not SFT, completing
the proof.
\end{proof}

\end{document}
